\section{Metodologia}



\subsection{Random Forest}
O Random Forest (RF) é um algoritmo de aprendizado de máquina supervisionado baseado em
conjuntos de árvores de decisão, no qual múltiplas árvores são treinadas a partir de
subconjuntos aleatórios dos dados e dos atributos, combinando suas previsões por meio de
votação majoritária. Sua utilização neste trabalho se deve à robustez frente a dados de
alta dimensionalidade, à capacidade de modelar relações não lineares e à possibilidade de
analisar a importância das features, características relevantes para problemas de
detecção de intrusão em redes~\cite{breiman_random_forest}.

\subsection{K-Nearest Neighbors (KNN)}
Como modelo classificador basal e de caráter não paramétrico, utilizou-se o algoritmo $K$-Nearest Neighbors 
(KNN). O KNN diferencia-se das abordagens estatísticas tradicionais por não construir um modelo matemático 
explícito durante a fase de treinamento, baseando sua inferência na busca por instâncias conhecidas armazenadas 
em memória que apresentem maior proximidade geométrica em relação à nova entrada~\cite{cover_hart_knn}.

Conforme detalhado na secção de pré-processamento, a aplicação da codificação binária (\textit{One-Hot Encoding}) 
e da padronização (\textit{StandardScaler}) foi de extrema importância para garantir a coerência matemática deste 
classificador. Sob a ótica do KNN, tais transformações asseguram que as variáveis categóricas não introduzam 
relações de ordem artificiais e que atributos com ordens de grandeza discrepantes assumam pesos idênticos no 
cálculo das distâncias espaciais.

Para a definição da configuração ideal, aplicou-se a estratégia de otimização automatizada por busca em grade 
(\textit{GridSearchCV}) detalhada anteriormente. O espaço de busca para este modelo englobou diferentes números 
de vizinhos ($K \in \{3, 5, 7, 9, 11, 15\}$) e métricas de distância (Euclidiana e Manhattan). O processo identificou 
como ótima a combinação de $K=3$ com a distância de Manhattan (norma $L_1$). A adoção de $K=3$ mitiga o risco de 
\textit{overfitting} inerente a $K=1$, enquanto a métrica de Manhattan demonstra-se teoricamente mais estável e eficiente 
do que a Euclidiana (norma $L_2$) para lidar com a Maldição da Dimensionalidade num espaço expandido de 118 \textit{features}.

\subsection{Multilayer Perceptron (MLP)}
O Multilayer Perceptron (MLP) é um modelo de rede neural artificial supervisionado,
composto por camadas de neurônios interconectados que aplicam transformações não lineares
sobre os dados de entrada. A escolha do MLP foi motivada por sua capacidade de aprender
padrões complexos no tráfego de rede, sendo adequado para a detecção de intrusões que
apresentam comportamentos sutis. 

Neste trabalho, o modelo foi avaliado em dois cenários:
um utilizando todas as features disponíveis e outro com a remoção da feature mais
influente identificada no modelo de melhor resultado, com o objetivo de analisar o impacto da
redução de atributos em seu desempenho. 

Ressalta-se que o desempenho do MLP depende
fortemente da escolha adequada de hiperparâmetros e do pré-processamento dos dados,
conforme discutido na literatura~\cite{haykin_neural_networks}.

\subsection{Regressão Logística}
A Regressão Logística (LR) é um algoritmo de aprendizado de máquina supervisionado clássico,
amplamente utilizado para a estimativa de probabilidades em problemas de classificação binária.
Diferentemente de abordagens baseadas em partições de entropia, o modelo opera construindo uma fronteira
de decisão linear através de uma combinação ponderada dos atributos de entrada ($w \cdot x$),
mapeando o valor contínuo resultante no intervalo estrito de $[0, 1]$ através da função logística Sigmoide.

A introdução da Regressão Logística neste trabalho justifica-se pela sua extrema eficiência computacional,
transparência paramétrica e por estabelecer uma linha de base estatística comparativa para os classificadores mais complexos.
Assim como no caso do KNN e do MLP, a normalização prévia das variáveis através da padronização de escala
(\textit{StandardScaler}) e a codificação binária de atributos nominais foram etapas obrigatórias. Tal tratamento
assegura que variáveis de elevada magnitude, como a contagem de bytes transferidos, não distorçam o gradiente de
otimização ou o cálculo dos coeficientes globais.

Alinhando-se estritamente à metodologia de avaliação dos demais modelos supervisionados, a Regressão Logística foi
testada sob dois regimes experimentais: o cenário completo, englobando a totalidade das variáveis preditoras, e o
cenário reduzido, no qual a \textit{feature} de maior impacto (\textit{src\_bytes}) foi suprimida. Esta abordagem permitiu
avaliar a resiliência estrutural e a capacidade de redistribuição automática de pesos do classificador linear perante
condições de perda de telemetria na rede.